\section{Directed-cycle consequences}\label{sec:directed-cycle-consequences}

In an oriented graph with no directed triangle, every exact second
out-neighbor is a non-neighbor.  Combining this observation with
Theorem~\ref{thm:main} gives degree and order bounds.
The first assertion below is the oriented-graph form of Thomass\'e's
non-neighbor statement recorded by
Sullivan~\cite[Conjecture~6.12]{Sullivan2006}.
The consequence using both minimum indegree and minimum outdegree is also
noted in \cite[Section~3]{Sullivan2006}.  The regular case gives the
girth-four lower bound in the Behzad--Chartrand--Wall minimum-order
conjecture~\cite{BehzadChartrandWall1970}, also recorded by
Sullivan~\cite[Section~4]{Sullivan2006}.

\begin{corollary}[Non-neighbors and directed cycles]\label{cor:directed-cycle-consequences}
Let \(D\) be a nonempty finite oriented graph of order \(n\).  For
\(v\in V(D)\), write
\[
 N_1^-(v)=\{u:u\to v\},\qquad
 d^\pm(v)=\abs{N_1^\pm(v)},\qquad
 \delta^\pm(D)=\min_{x\in V(D)}d^\pm(x).
\]
Define the non-neighbors of \(v\), excluding \(v\) itself, by
\[
 Z(v)=V(D)\setminus
 \bigl(\{v\}\cup N_1^+(v)\cup N_1^-(v)\bigr).
\]
Then the following hold.
\begin{enumerate}[label=\textup{(\roman*)}]
\item If \(D\) has no directed \(3\)-cycle, then some vertex \(v\) satisfies
\[
 d^+(v)\leq\abs{Z(v)},\qquad
 n-1\geq 2d^+(v)+d^-(v).
\]
\item If \(\delta^+(D)\geq n/3\) and \(\delta^-(D)\geq n/3\), then \(D\)
contains a directed \(3\)-cycle.
\item If \(r\geq1\) is an integer, \(d^+(v)=d^-(v)=r\) for every vertex,
and \(D\) has directed girth exactly four, then \(n\geq3r+1\).
\end{enumerate}
Here directed girth is the length of a shortest directed cycle.
\end{corollary}

\begin{proof}
Suppose that \(D\) has no directed \(3\)-cycle.  If \(w\in N_2^+(v)\),
choose a path \(v\to u\to w\).  The exact-distance definition excludes
\(w=v\) and \(w\in N_1^+(v)\).  If \(w\in N_1^-(v)\), the arc \(w\to v\)
would close a directed \(3\)-cycle.  Thus
\(N_2^+(v)\subseteq Z(v)\) for every \(v\).  Theorem~\ref{thm:main}
now gives a vertex with
\[
 d^+(v)\leq\abs{N_2^+(v)}\leq\abs{Z(v)}.
\]
Since \(D\) is oriented, its first in- and out-neighborhoods are disjoint,
and their cardinalities are the corresponding degrees.  Hence
\[
 n-1=d^+(v)+d^-(v)+\abs{Z(v)}
 \geq 2d^+(v)+d^-(v),
\]
proving \textup{(i)}.

Under the hypotheses of \textup{(ii)}, absence of a directed \(3\)-cycle
would therefore give
\[
 n-1\geq2\delta^+(D)+\delta^-(D)\geq n,
\]
a contradiction.  The integer degree thresholds are
\(\delta^\pm(D)\geq\lceil n/3\rceil\), so this includes the endpoint
when \(3\) divides \(n\).  Finally, directed girth four excludes directed
\(3\)-cycles; in the regular case the count in \textup{(i)} becomes
\(n-1\geq3r\), proving \textup{(iii)}.
\end{proof}

